\section*{Chapter \ref{ch:qm:numerical-optimisation-in-practice} --- Numerical Optimisation in Practice}

\begin{solution}{exo:qm:numerical-optimisation-in-practice:1}
With the bound $(1 - 1/\kappa)^k$, $k \approx \kappa\ln10^6 = 13\,816$ iterations; with acceleration, about $\sqrt\kappa\ln10^6 = 437$.
\end{solution}

\begin{solution}{exo:qm:numerical-optimisation-in-practice:2}
Soft thresholding by 1: $(2, 0, 0)$.
\end{solution}

\begin{solution}{exo:qm:numerical-optimisation-in-practice:3}
Swapping $(a, \tau_1, \tau_2)$ for $(1 - a, \tau_2, \tau_1)$ gives the same function of the lag: the model is invariant under relabelling its components, so every minimum has a mirror image.
Ordering the time scales in the parameterisation removes it.
\end{solution}

\begin{solution}{exo:qm:numerical-optimisation-in-practice:4}
For $f(x) = \frac12x^\top Hx + b^\top x$ with $H$ positive definite, $x - H^{-1}(Hx + b) = -H^{-1}b$, the unique minimiser, whatever $x$.
\end{solution}

\begin{solution}{exo:qm:numerical-optimisation-in-practice:5}
With $s = td$, $y = \nabla f(x + td) - \nabla f(x)$: $y^\top s = t(\nabla f(x + td)^\top d - \nabla f(x)^\top d) \ge t(c_2 - 1)\nabla f(x)^\top d > 0$, since $\nabla f(x)^\top d < 0$ and $c_2 < 1$.
\end{solution}

\begin{solution}{exo:qm:numerical-optimisation-in-practice:6}
$1/\sqrt k$: the sum diverges but so does the sum of squares ($\sum1/k$), so not in the classical form (though averaging the iterates rescues it). $1/k$: the sum diverges and the sum of squares
converges, so yes.
\end{solution}

\begin{solution}{exo:qm:numerical-optimisation-in-practice:7}
\texttt{daily(w)} for $w = 0$, 0.001, 0.01, 0.1: median change 4.48, 3.93, 1.87 and 0.42 days; average RMSE ($\times 10^{-3}$) 9.450, 9.454, 9.520 and 9.643. Stability is cheap up to 0.01 and then starts to
cost fit.
\end{solution}

\begin{solution}{exo:qm:numerical-optimisation-in-practice:8}
Convergence to a stationary point says nothing about which minimum was found (mirror images, boundary solutions), nor about how flat the objective is there: parameters in a flat valley move
with noise while the gradient is zero at each day's point. Check \omterm{def:qm:numerical-optimisation-in-practice:calibration}{multistart}, the condition number of $J^\top J$ and the day-to-day stability.
\end{solution}

\begin{solution}{pb:qm:numerical-optimisation-in-practice:1}
\textbf{1.} $a = 0.59$, $\tau_1 = 2.9$ and $\tau_2 = 55$ days, RMSE 0.0096.
\textbf{2.} 28 at the best fit, 35 at its mirror image, 26 at the 1\,000-day boundary with at least twice the error, 11 elsewhere.
\textbf{3.} $(0.41, 55, 2.9)$: the same curve with the components' labels swapped.
\textbf{4.} Condition number 10.5; correlation 0.96 between the fast weight and the slow time scale.
\textbf{5.} From about 45 to 80 days (profile error $9.9 \times 10^{-3}$ at 62 days, 12.3 at 42, 12.9 at 89, against noise of 10).
\textbf{6.} Median 4.5 days, maximum 30.7.
\textbf{7.} Day 189: the slow time scale went from 56.8 to 87.6 days and the fit error from 0.00926 to 0.00921.
\textbf{8.} Median 1.9 days, maximum 9.5.
\textbf{9.} The average RMSE rises from 0.00945 to 0.00952, 0.7\%.
\textbf{10.} Its standard deviation falls from 5.5 to 2.7 days.
\textbf{11.} \omterm{def:qm:numerical-optimisation-in-practice:gd}{Gradient descent} reaches 0.067 after 100 iterations and 0.0090 after 200; Nesterov $5 \times 10^{-6}$ after 100.
\textbf{12.} One Newton step (two iterations including the check); L-BFGS seven.
\textbf{13.} ISTA is $7.6 \times 10^{-4}$ from the optimum after 50 iterations; FISTA about $10^{-7}$ after 50 and below $10^{-10}$ after 65.
\textbf{14.} 91 ADMM iterations against 74 projections, the same matrix to $1.4 \times 10^{-9}$.
\textbf{15.} Convergence: 0.009 from the least-squares solution after 20\,000 steps, against 0.113 with a constant step.
\textbf{16.} The slow time scale, with its interval along the valley (about 45 to 80 days) and its day-to-day change.
\textbf{17.} By \omterm{def:qm:linear-models-under-stress:cv}{cross-validation}: the weight that minimises tomorrow's fit error of today's parameters, or a weight that caps the day-to-day change at a tolerance the users accept.
\textbf{18.} Parameterise $\tau_2 = \tau_1 + e^u$ (and $a \in (0, 1)$), so only one labelling is representable.
\textbf{19.} Named result: \emph{the Monday the fit moved}: refit daily from a fixed start, the slow time scale jumps by a median of 4.5 days (maximum 30.7); with a penalty toward yesterday (weight
0.01) by 1.9 (maximum 9.5), at a cost of 0.7\% in fit error.
\textbf{20.} Which of several equally good parameter sets it found, and how precisely the data determine them.
\end{solution}

\begin{solution}{iq:qm:numerical-optimisation-in-practice:1}
The step must be small enough for the stiffest direction ($1/L$), so progress along the flattest direction is only $\mu/L$ per step: the rate is set by the condition number. Remedies: rescale the
variables, precondition, use curvature (Newton, quasi-Newton, Gauss--Newton), or accelerate (Nesterov, momentum).
\iqlookfor{Condition number as the rate, and at least two remedies.}
\end{solution}

\begin{solution}{iq:qm:numerical-optimisation-in-practice:2}
The objective is flat in some direction (an \omterm{def:qm:numerical-optimisation-in-practice:calibration}{identifiability} problem), or has several equivalent minima, and noise moves the answer. Diagnose with the condition number of $J^\top J$, a profile of
the objective and \omterm{def:qm:numerical-optimisation-in-practice:calibration}{multistart}; fix with a penalty toward yesterday, reparameterisation, fixing a poorly identified parameter, or more informative data.
\iqlookfor{\omterm{def:qm:numerical-optimisation-in-practice:calibration}{Identifiability} diagnosis, and a stabilisation that costs little fit.}
\end{solution}

\begin{solution}{iq:qm:numerical-optimisation-in-practice:3}
It solves $(J^\top J + \lambda D)\Delta x = -J^\top r$: small $\lambda$ gives Gauss--Newton (fast near the solution), large $\lambda$ a short scaled gradient step (safe far from it); $\lambda$ adapts to
success. Plain Gauss--Newton can take huge steps when $J^\top J$ is nearly singular, exactly in flat valleys.
\iqlookfor{Interpolation between Gauss--Newton and \omterm{def:qm:numerical-optimisation-in-practice:gd}{gradient descent}, and the singular $J^\top J$.}
\end{solution}

\begin{solution}{iq:qm:numerical-optimisation-in-practice:4}
Each stochastic gradient has variance that does not vanish at the optimum, so a constant step keeps the iterate bouncing in a noise ball of radius proportional to the step. With $\sum t_k = \infty$
(enough total movement) and $\sum t_k^2 < \infty$ (vanishing noise), the iterates converge.
\iqlookfor{Noise at the optimum and the two conditions.}
\end{solution}

\begin{solution}{iq:qm:numerical-optimisation-in-practice:5}
$\operatorname{prox}_{tg}(v)$ minimises $g(x) + \lVert x - v\rVert^2/2t$; for the $\ell_1$ norm it is soft thresholding. The \omterm{def:qm:linear-models-under-stress:penalties}{lasso}'s non-smooth penalty is handled exactly by the prox, while the smooth
least-squares part uses its gradient: the \omterm{def:qm:numerical-optimisation-in-practice:prox}{proximal gradient method} (ISTA, FISTA) and coordinate descent both exploit it.
\iqlookfor{Soft thresholding and the smooth-plus-simple split.}
\end{solution}

\begin{solution}{iq:qm:numerical-optimisation-in-practice:6}
The BFGS update keeps the inverse-Hessian approximation positive definite only if $y^\top s > 0$; the Wolfe curvature condition guarantees it, while a mere decrease condition does not. Without it
the next direction may not descend.
\iqlookfor{Curvature condition implies positive definiteness.}
\end{solution}
